% SOURCE_AUTHORITY: sga5_fr_workpass.tex, lines 9171--9199 (LF-normalized unit).
% SOURCE_SHA256: 791F4EFFC5E02832D5D77ED03518C8156D6F07E4C8238B03545DB93D883FBB28
\emph{Demonstração.} Sejam $\alpha:U\to X$ e $\beta:Y\to X$ as canônicas imersões. Sabe-se (SGA 4 XVI 3.7) que, definindo $F=p^*(L)$, é
\[
F\xrightarrow{\sim}\alpha_*\alpha^*(F)
\quad\text{e}\quad
(\R^j\alpha_*)\alpha^*(F)=0\quad\text{se }j\ne0,2r-1.
\]
Por outro lado, os feixes $(\R^j\alpha_*)(\alpha^*F)$, para $j\ge1$, têm suporte em $Y$ e são, portanto, acíclicos para o funtor $p_*$. A partir destas observações, deduzem-se as afirmações (i) a) e (b) utilizando a sucessão espectral de Leray
\[
\R^i p_*\bigl(\R^j\alpha_*\bigr)(q^*L)\Longrightarrow \R^n q_*(q^*L).
\]
A afirmação (i) c) segue-se de (1.1 (i)) e de longa sequência exata
\[
\cdots\longrightarrow \R_Y^i p_*(p^*L)\longrightarrow \R^i p_*(p^*L)\longrightarrow \R^i q_*(q^*L)
\longrightarrow \R_Y^{i+1}p_*(p^*L)\longrightarrow\cdots.
\]
O feixe $\R_Y^{2r}p_*(p^*L)$ é associado ao prehaz
\[
U\longmapsto H^{2r}_{Y\times_S U}(X\times_S U,p^*L).
\]

Demostremos (ii). A afirmação c) segue-se de que o morfismo $q$ é suave e tem fibras geométricas relacionadas de dimensão $r$. A sequência exata mencionada em (b) é obtida aplicando o funtor $\R_!p$ à canónica sequência exata
\[
0\longrightarrow \alpha_!(q^*L)\longrightarrow p^*L\longrightarrow \beta_!(\beta^*p^*L)\longrightarrow 0.
\]
A utilização de (1.1 (i) permite deduzir sem dificuldade das afirmações (ii) a) e (b) deste sequência exata longo.

\begin{statement}{Observação 1.4}
Para $S=\operatorname{spec}(\mathbb C)$, (1.3) expressa que o complementário de origem em $\mathbb C^r$ tem a mesma cohomologia que a esfera real de dimensão $2r-1$, retrata-se por deformação de $\mathbb C^r-0$.
\end{statement}
